\section{Introduction}
\label{sec:intro}

Deep learning models have experienced rapid advancements in various areas, such as 2D vision~\cite{tian2021divide,goyal2021self,wang2022painter,kirillov2023segany} and natural language processing (NLP)~\cite{kaplan2020scaling,aribandi2022ext,touvron2023llama,openai2023gpt4}, with their progress often attributed to the effective utilization of scale, encompassing factors such as the size of datasets, the number of model parameters, the range of effective receptive field, and the computing power allocated for training.

However, in contrast to the progress made in 2D vision or NLP, the development of 3D backbones~\cite{qi2017pointnet, dai20183dmv, li2018pointcnn, wu2019pointconv} has been hindered in terms of scale, primarily due to the limited size and diversity of point cloud data available in separate domains~\cite{wu2023ppt}. Consequently, there exists a gap in applying scaling principles that have driven advancements in other fields~\cite{kaplan2020scaling}. This absence of scale often leads to a limited trade-off between accuracy and speed on 3D backbones, particularly for models based on the transformer architecture~\cite{zhao2021point, guo2021pct}. Typically, this trade-off involves sacrificing efficiency for accuracy. Such limited efficiency impedes some of these models' capacity to fully leverage the inherent strength of transformers in scaling the range of receptive fields, hindering their full potential in 3D data processing. \looseness=-1

A recent advancement~\cite{wu2023ppt} in 3D representation learning has made progress in overcoming the data scale limitation in point cloud processing by introducing a synergistic training approach spanning multiple 3D datasets. Coupled with this strategy, the efficient convolutional backbone~\cite{choy20194d} has effectively bridged the accuracy gap commonly associated with point cloud transformers~\cite{wu2022point, lai2022stratified}. However, point cloud transformers themselves have not yet fully benefited from this privilege of scale due to their efficiency gap compared to sparse convolution. This discovery shapes the initial motivation for our work:  \textit{to re-weigh the design choices in point transformers, with the lens of the scaling principle}. We posit that model performance is more significantly influenced by scale than by intricate design.

Therefore, we introduce Point Transformer V3 (PTv3), which prioritizes simplicity and efficiency over the accuracy of certain mechanisms, thereby enabling scalability. Such adjustments have an ignorable impact on overall performance after scaling. Specifically, PTv3 makes the following adaptations to achieve superior efficiency and scalability:\looseness=-1
\begin{itemize}[leftmargin=4mm, itemsep=0mm, topsep=0mm, partopsep=0mm]
    \item Inspired by two recent advancements~\cite{Wang2023OctFormer, liu2023flatformer} and recognizing the scalability benefits of structuring unstructured point clouds, PTv3 shifts from the traditional spatial proximity defined by K-Nearest Neighbors (KNN) query, accounting for 28\% of the forward time. Instead, it explores the potential of serialized neighborhoods in point clouds, organized according to specific patterns.
    \item PTv3 replaces more complex attention patch interaction mechanisms, like shift-window (impeding the fusion of attention operators) and the neighborhood mechanism (causing high memory consumption), with a streamlined approach tailored for serialized point clouds.
    \item PTv3 eliminates the reliance on relative positional encoding, which accounts for 26\% of the forward time, in favor of a simpler prepositive sparse convolutional layer.
\end{itemize}
We consider these designs as intuitive choices driven by the scaling principles and advancements in existing point cloud transformers. Importantly, this paper underscores the critical importance of recognizing how scalability affects backbone design, instead of detailed module designs.

This principle significantly enhances scalability, overcoming traditional trade-offs between accuracy and efficiency (see \figref{fig:teaser}). PTv3, compared to its predecessor, has achieved a 3.3$\times$ increase in inference speed and a 10.2$\times$ reduction in memory usage. More importantly, PTv3 capitalizes on its inherent ability to scale the range of perception, expanding its receptive field from 16 to 1024 points while maintaining efficiency. This scalability underpins its superior performance in real-world perception tasks, where PTv3 achieves state-of-the-art results across over 20 downstream tasks in both indoor and outdoor scenarios. Further augmenting its data scale with multi-dataset training~\cite{wu2023ppt}, PTv3 elevates these results even more. We hope that our insights will inspire future research in this direction.
